% Source: https://artofproblemsolving.com/community/c3h2239606_axbyc_and_arthemtic_progression
% Stable identity: aops-topic-2239606-q01
The point $(1,-2)$ lies on the line exactly when
\[A\cdot1+B\cdot(-2)+C=0,\qquad\text{that is}\qquad A+C=2B .\]
And $A$, $B$, $C$ form an arithmetic progression exactly when $B-A=C-B$, which is the same condition $A+C=2B$.

\textbf{(a)} If $A,B,C$ is an arithmetic progression, then $A+C=2B$, so $A-2B+C=0$ and $(1,-2)$ is on the line.

\textbf{(b)} If $(1,-2)$ is on the line, then $A-2B+C=0$, so $B-A=C-B$ and $A,B,C$ is an arithmetic progression. $\blacksquare$

For example, $(A,B,C)=(1,1,1)$, $(1,2,3)$, $(2,1,0)$ and $(3,1,-1)$ give four lines through $(1,-2)$:

\begin{asy}
usepackage("amsmath");
size(6cm);
defaultpen(fontsize(10pt));
pen[] pp={blue,red,heavygreen,orange};
real[][] L={{1,1,1},{1,2,3},{2,1,0},{3,1,-1}};
string[] nm={"$x+y+1=0$","$x+2y+3=0$","$2x+y=0$","$3x+y-1=0$"};
for(int i=0;i<4;++i){
  real a=L[i][0], bq=L[i][1], c=L[i][2];
  draw((-3,(-c+3a)/bq)--(4,(-c-4a)/bq),pp[i]);
}
clip(box((-2.5,-4.5),(3.6,2.6)));
draw((-2.5,0)--(3.6,0),Arrow(6)); draw((0,-4.5)--(0,2.6),Arrow(6));
label("$x$",(3.6,0),S); label("$y$",(0,2.6),W);
for(int i=0;i<4;++i){
  real yy=2.2-0.6i;
  draw((0.8,yy)--(1.3,yy),pp[i]); label(nm[i],(1.3,yy),E,pp[i]+fontsize(8pt));
}
dot((1,-2)); label("$(1,-2)$",(1,-2),E,fontsize(8pt));
\end{asy}
